\section{矩阵补全从缺失条目到低秩正则}
\label{sec:13-Machine-Learning/04-Dimensionality-Reduction/03}

手上是一张用户--物品评分表，几百万行、几十万列，填了数字的格子不到千分之三。产品要你把空格补上。直觉很顺：用户的口味类型有限，物品的属性维度也有限，几十个隐向量就能把两边都概括掉，整张表应该近似低秩；找到那个低秩结构，空白自然浮现。这个直觉对了一大半。剩下那一小半——低秩把候选压缩到多小、观测落在哪里、以及为什么某些格子无论如何都补不出来——是本节要说清楚的。

\subsection{低秩声明的是自由度很少，不是答案唯一}

前两节的降维都在丢方向：PCA 丢方差小的方向，截断 SVD 丢强度弱的通道。矩阵补全反过来，它先声明``这张表只有少数几个自由度''，再问在这个声明之下，观测到的条目能不能把整张表钉死。

答案是经常不能。取一个 \(2\times2\) 的例子，只观测对角线，左上是 1、右下是 4。下面两个矩阵

\[
\begin{pmatrix}1&2\\2&4\end{pmatrix}
\qquad\text{与}\qquad
\begin{pmatrix}1&1\\4&4\end{pmatrix}
\]

都是秩一，对角线完全相同，非对角线一个是 2、一个分别是 1 和 4。秩约束把候选集从整个 \(\R^{2\times 2}\) 压到了一张二维曲面，可它没有压到一个点。

更极端的情形是某个用户一条评分都没有。他那一整行可以填任何数——只要与其余行拼起来仍然低秩就行，而这总能做到。这时候再换更大的模型、跑更多轮迭代、把 \(\lambda\) 调得更细，都不会有帮助：信息不在数据里，算法变不出来。低秩是一个关于整张表的全局约束，它不生产局部证据。

\subsection{秩换成核范数，是为了把不可解的问题换成可解的}

理想形式很干净。记真实矩阵为 \(M\)，\(\Omega\) 是观测位置集合，\(P_\Omega\) 保留这些位置的条目、其余置零，我们想要与观测一致的最低秩矩阵：

\[
\min_X\ \rank(X)
\quad\text{s.t.}\quad
P_\Omega(X)=P_\Omega(M).
\]

干净但不可解。\(\rank\) 取整数值、随矩阵跳变，这类问题一般是 NP-hard 的。处理办法和稀疏回归里的一模一样：向量里用 \(\ell_1\) 替代 \(\ell_0\)，矩阵里用核范数替代秩，

\[
\norm{X}_* = \sum_j\sigma_j(X),
\]

也就是奇异值向量的 \(\ell_1\) 范数。Lasso 把系数往零压，核范数把奇异值往零压，而每多一个零奇异值，秩就掉一。替换之后问题变凸。带噪声时更常写成无约束形式：记 \(Y\) 为观测到的带噪条目，

\[
\min_X\ \tfrac12\norm{P_\Omega(X-Y)}_F^2 + \lambda\norm{X}_* .
\]

第一项忠于观测，第二项压低复杂度，\(\lambda\) 决定两者的兑换率。解它的核心操作是奇异值软阈值：对当前矩阵做 SVD，把每个 \(\sigma_j\) 换成 \((\sigma_j-\tau)_+\)，再乘回去。迭代软阈值算法就是在``沿梯度走半步''和``过一遍软阈值''之间交替。

另一条路线把矩阵显式写成两个瘦因子的乘积 \(X=UV^{\trans}\)，目标改成

\[
\tfrac12\norm{P_\Omega(UV^{\trans}-Y)}_F^2
+\tfrac{\lambda}{2}\bigl(\norm{U}_F^2+\norm{V}_F^2\bigr).
\]

这个形式不必存下整个 \(n\times m\) 矩阵，大规模场景下几乎是唯一可行的做法。要认下的麻烦有三条：目标对 \((U,V)\) 不再凸，秩必须事先选定，以及因子本身不唯一——对任意可逆 \(R\)，\((UR,\,VR^{-\trans})\) 给出同一个乘积。最后一条经常被忽略，后果是有人指着 \(U\) 的某一列讲``这维代表消费能力''。那是讲不通的：把 \(U\) 旋转一下，这列的含义就变了，而乘积一个数都没动。可解释的是列空间，不是列。

\subsection{观测掩码一出现，截断的最优性就失效了}

值得停下来对一次账。上一节的 Eckart--Young 保证截断 SVD 是最优低秩近似，条件是可行集恰好是 \(\set{B:\rank(B)\le k}\)、范数酉不变、没有别的限制。这里的 \(P_\Omega\) 恰恰是别的限制：误差只在观测位置上计量，\(\norm{P_\Omega(\cdot)}_F\) 不是酉不变范数——左右乘一个正交阵，掩码就对不上了。

所以``把缺失位置填零、然后做截断 SVD''这个看起来自然的做法没有最优性保护，而且偏差方向可以预判：填零等于断言未观测的评分是 0，低秩拟合会被这批假零往下拽。用行均值或列均值填充也只是把偏差换个方向。矩阵补全之所以要写成一个优化问题反复迭代，正是因为它不能像上一节那样一步分解出答案——这是那条边界在实际问题里的第一次兑现。

\subsection{可恢复性要求结构与采样互相配合}

凸松弛本身不承诺恢复。精确恢复的定理需要几件事同时成立：矩阵确实足够低秩；奇异向量不过度集中在少数坐标上，即所谓 incoherence；观测数量足够且在行列间分布得开；噪声受控。

incoherence 为什么关键，一个例子就够了。设想一个秩一矩阵，唯一的非零条目在左上角，其余全是零。除非你恰好采到左上角那格，你看到的数据与全零矩阵毫无区别，而秩零的解已经完美匹配观测。低秩说的是自由度少，incoherence 说的是这些自由度没有藏在极个别条目里。信息越集中，需要的采样越密。

\begin{center}
\begin{tikzpicture}[>=stealth,scale=1.0]
% ---- 左：分散的观测 ----
\foreach \i in {0,1,2,3,4,5,6,7}{
  \draw[gray!35] ({\i*0.34},0) -- ({\i*0.34},-2.38);
  \draw[gray!35] (0,{-\i*0.34}) -- (2.38,{-\i*0.34});
}
\foreach \c/\r in {0/0,2/0,4/0,6/0, 1/1,3/1,5/1, 0/2,2/2,3/2,6/2,
                   1/3,4/3,5/3, 0/4,2/4,5/4,6/4, 1/5,3/5,4/5, 0/6,2/6,4/6,6/6}
  \fill[black!72] ({\c*0.34+0.05},{-\r*0.34-0.05}) rectangle ({\c*0.34+0.29},{-\r*0.34-0.29});
\node[font=\small] at (1.19,-2.85) {观测落得开};

% ---- 右：集中的观测 + 空行 ----
\begin{scope}[shift={(4.6,0)}]
\foreach \i in {0,1,2,3,4,5,6,7}{
  \draw[gray!35] ({\i*0.34},0) -- ({\i*0.34},-2.38);
  \draw[gray!35] (0,{-\i*0.34}) -- (2.38,{-\i*0.34});
}
\foreach \c/\r in {0/0,1/0,2/0, 0/1,1/1,2/1, 0/2,1/2,2/2,5/2,
                   0/4,1/4, 0/5,1/5,2/5, 0/6,1/6,2/6,6/6}
  \fill[black!72] ({\c*0.34+0.05},{-\r*0.34-0.05}) rectangle ({\c*0.34+0.29},{-\r*0.34-0.29});
\draw[black,->] (3.05,-1.19) -- (2.48,-1.19);
\node[font=\scriptsize,anchor=west] at (3.1,-1.19) {整行为空};
\node[font=\small] at (1.19,-2.85) {观测挤在少数列};
\end{scope}
\end{tikzpicture}
\end{center}

\emph{同样是低秩矩阵、观测数量也相近，左边每一行每一列都拿到证据，右边把预算耗在头几列上，还有一行完全没有采样。}

图右侧那一行的空白说明了一件事：可恢复性不是矩阵的属性，也不是算法的属性，而是矩阵结构与采样模式的联合属性。换个掩码，同一个矩阵可能从能恢复变成不能恢复。

推荐场景里的麻烦比 incoherence 更深一层：缺失几乎从来不是随机的。用户倾向于给特别喜欢或特别厌恶的内容打分，平台又优先展示算法猜他会喜欢的东西。``没有评分''可能意味着没兴趣，也可能意味着根本没被展示过，这两种含义在数据里长得一模一样。未观测本身携带信号，而上面那个优化问题默认它不携带。低秩假设可以在给定缺失机制下预测未观测条目，它没有能力把``选择性曝光下的评分''翻译成``随机曝光下的真实偏好''，那是一个因果问题，需要曝光概率或随机化实验才能回答。

\subsection{验证方式必须复制真正的缺失方式}

最常见的陷阱是把已观测评分随机拆成训练和测试。这样同一个用户和同一个物品在两边大量出现，测出来的是``熟悉用户对熟悉物品''的插值能力，而上线之后你面对的是新用户、新物品、以及下个月的分布。两个数字量的不是一回事，前者通常好看得多。

正确的切法取决于部署情境：按用户留出评估冷启动，按物品留出评估新品上架，按时间切分评估分布漂移。冷启动用户一条训练评分都没有，只能靠其他用户学到的物品结构去推；新物品则反过来。这几种切法测的是不同的泛化能力，报告时应当分开，而不是合成一个总分。这条要求与\hyperref[sec:11-Mathematical-Statistics/06-Resampling-and-Model-Selection/03]{``交叉验证评估的是完整训练流程''}一节说的是同一件事：被评估的是整条流程连同它的假设，不只是末端那个模型。

另外几件事值得成为固定动作。先把观测掩码画出来，看清整行整列的空缺和曝光偏差；再和最朴素的基线比一比——用户平均分加物品平均偏移，如果低秩模型打不过它，说明学到的不是结构而是噪声；最后用合成的低秩矩阵做一次已知答案的恢复实验，然后只改掩码、不改矩阵，看误差怎么变。这三步能挡掉大部分自欺。

有时候最诚实的结论是：这一行确实没有信息，当前这张表识别不了它，需要侧信息或主动采集。不过``低秩''本身也可能是错的假设——如果用户偏好沿着一条弯曲的结构分布，任何线性因子模型都会系统性地偏。下一节把这个假设换掉。
